\documentclass[border=4pt]{standalone}
\usepackage{amsmath,amssymb,mathtools,xcolor,varwidth}
\definecolor{arearred}{HTML}{D81B60}
\begin{document}
\begin{varwidth}{28em}
\large\bfseries
Now hold $Ae$ fixed and let $B$ and $C$ meet $A$; the angle $cAg$ between curve      and tangent vanishes, exactly as in Lemma VII, so the curvilinear areas $Abd$ and $Ace$      coincide with the \textcolor{arearred}{rectilinear triangles $Afd$ and $Age$}. But these two triangles      share the same vertex angle at $A$ along the tangent $Ag$, so by \textcolor{arearred}{Lemma V (similar      rectilinear figures)} their areas are as the squares of the homologous sides, that is      $Afd : Age = Ad^2 : Ae^2$. Since $Ad : Ae = AD : AE$ and the original triangles $\triangle      ABD$, $\triangle ACE$ stay always proportional to these blown-up areas, the same duplicate      ratio descends to them. Therefore the areas $\triangle ABD$ and $\triangle ACE$ are      ultimately one to the other as $AD^2 : AE^2$. \textit{Q.E.D.}
\end{varwidth}
\end{document}
